\chapter{Myelography}
\section*{What Is This Test?} Myelography uses contrast and X-ray or CT imaging to show the spinal cord, nerve roots, and spinal canal.\section*{Alternative Names} CT myelogram; intrathecal contrast study; myelogram.\section*{Principle} Contrast injected into cerebrospinal fluid outlines spaces that may be narrowed or displaced by disc disease, stenosis, tumor, or other lesions.\section*{Why This Test Is Ordered} It evaluates nerve compression, spinal stenosis, tumors, or structural disease when MRI is unavailable, unsuitable, or inconclusive.\section*{Primary vs Secondary Test} It is a secondary invasive imaging study; MRI is often preferred when feasible.\section*{How This Test Is Performed} A needle enters the spinal canal under imaging guidance, contrast is injected, and radiographs or CT images are acquired.\section*{What Preparation Is Needed} Fasting and medicine instructions vary. Report pregnancy, bleeding disorders, kidney disease, seizures, and contrast reactions.\section*{Understanding Results} Findings may show narrowing, blockage, nerve-root displacement, or mass effect; interpretation requires symptoms and examination.\section*{Follow-up Tests} MRI, surgical or neurologic consultation, or additional imaging may follow.\section*{References} \begin{enumerate}\item MedlinePlus. Myelography.\item American College of Radiology. Myelography practice guidance.\end{enumerate}\clearpage\section*{Understanding Results and Follow-up}
The laboratory report should identify the specimen, method, units, reference interval or decision limit, and whether the result is positive, negative, abnormal, or inconclusive. Reference intervals vary with age, sex, pregnancy, specimen type, assay, and laboratory; a result outside a range is not automatically a diagnosis, and a result within a range does not always exclude disease. Discuss unexpected results with the ordering clinician, who can decide whether repeat or confirmatory testing, treatment, or referral is appropriate. Do not change medicines or delay urgent care based on an isolated result.
\section*{Regulatory Status and Authorization Date}This chapter describes a test category, not a specific commercial product. FDA, CDSCO, EU IVDR, and MHRA approval or registration dates therefore cannot be assigned without the assay manufacturer, product name, instrument, and intended use. For laboratory-developed tests, an individual agency approval date may not exist. See the Preface for the interpretation of regulatory status.
\clearpage
